\textbf{Desarrollo de las magnitudes intermedias.} La factibilidad por sustentación se contrasta con $C_{L,max}$; la de motor compara la resistencia de ese mismo estado con el empuje disponible. Ninguna de las dos verifica por sí sola el trim de momentos.\begin{enumerate}

\item Fuerza de referencia de la fase a TAS fija.

\[\fitmath{\mathrm{qS}=0.5\cdot 0.623843718\cdot {175}^{2}\cdot 52\approx 496735.6\,\mathrm{N}}\]

\item Sustentación requerida por la masa de enlace.

\[\fitmath{C_L=\frac{11075.13882\cdot 9.80665}{\mathrm{qS}}}\]

\[\fitmath{C_L=\frac{11075.13882\cdot 9.80665}{(496735.5604)}\approx 0.2186475\,}\]

\item Resistencia y empuje requerido.

\[\fitmath{D=\mathrm{qS}\cdot \left(0.024+0.044\cdot {C_L}^{2}\right)}\]

\[\fitmath{D=(496735.5604)\cdot \left(0.024+0.044\cdot {(0.2186475436)}^{2}\right)\approx 12966.54\,\mathrm{N}}\]

\item Margen local de empuje disponible.

\[\fitmath{\mathrm{margin}=10500-D}\]

\[\fitmath{\mathrm{margin}=10500-(12966.53518)\approx -2466.535\,\mathrm{N}}\]

\end{enumerate}
